\subsection{Problem 18. Improper integrals and comparison}
Consider the improper integral
\[J(p)=\displaystyle\int_{1}^{\infty}\frac{1}{x^{p}}dx, \quad p>0\]

\begin{enumerate}[label=(\alph*)]
    \item Evaluate $J(p)$ explicitly for all $p>1$, and show that it diverges for $0<p\le1$.
    \item Now consider the integral $\displaystyle\int_{1}^{\infty}\frac{\ln x}{x^{p}}dx$. Using comparison or limit comparison, determine for which values of $p$ this integral converges.
    \item For a specific value of $p$ where convergence holds (for example $p=3$), compute the integral $\displaystyle\int_{1}^{\infty}\frac{\ln x}{x^{3}}dx$ explicitly using integration by parts.
    \item Explain how these examples illustrate the idea that a slowly growing factor like $\ln x$ can still be dominated by a sufficiently fast decaying power $x^{-p}$.
\end{enumerate}

\subsection*{Strategy}
\subsection*{Solution}
\pagebreak